\section{魏的 dossier：把旧都变成可调用的中枢}
\label{sec:wei-foundation-regency}

\begin{event}{\eventanchor{TK-0220-WEI-LUOYANG}{魏定都洛阳}}{黄初元年（二二〇年）}{曹魏 · 洛阳}{定都可核，城内复用与规模未详}
《三国志·文帝纪》《三国志·明帝纪》与《资治通鉴》魏纪可以互校曹丕定都洛阳的年序。宫城如何复用、营建节奏和城中人口规模，则仍须由遗址分期检验，不能从定都诏令直接推得。
\end{event}

\subsection{纲要：北方资源与中枢的近距离}

黄初元年，曹丕把受禅后的朝廷置于洛阳：
\hyperlink{evt:TK-0220-WEI-LUOYANG}{魏定都洛阳}。这项选择把尚书、禁军和北方仓储
重新接到旧都及其道路上，也让新朝不必继续完全依赖许县的战时行营格局。它首先
是一种可调用军政资源的布置，而不只是借用东汉旧都的象征。关中、淮南、辽东和
汉中的兵事仍各受道路与季节限制，洛阳发出的命令到达边地，更不等于能够按原样
执行。《文帝纪》与《资治通鉴》可确定定都的年序；宫城如何复用、营建有多快、
城中人口有多少，则须由遗址分期核验。

二二六年曹丕去世，\eventanchor{TK-DENS-0226-01}{曹叡即位并以曹真、司马懿等分掌
内外军政}，说明魏初的继承已经能由近畿机构接续；
卒年与继位可以由帝纪互证，遗诏所划出的
辅政层级和各集团的边界却不能据此写实。曹叡身后留下幼帝与两位辅政者，争夺
遂从皇位本身转向谁能调度近畿军队、诏令和任官。辽东的
\eventanchor{TK-0238-SIMA-YI-LIAODONG}{远征胜利}扩大了中枢声望，却没有替幼主
解决这道辅政接口；襄平围城后的处置与治理速度，也不能由征讨的结果倒推为
东北已经稳定整合。

\subsection{选官、营建与双辅政的缝隙}

受禅之年采纳的九品官人法，把察举已受战乱
打断的难题转成中正对士人资望的品评；它给洛阳一个安排官员的语言，也把地方
关系带入选官。制度可从《三国志·魏书》与《资治通鉴》互见，却不能据条文
断言各州已用同一种尺度。青龙三年，明帝继续营建宫殿、园囿和土木工程，朝臣
对役使提出异议。
记载足以把宫门内的工程和劳力压力并置，不能补出徭役人数；但它说明北方中枢
的可调用资源并非只流向边防。

景初三年，曹叡以曹爽、司马懿同受辅政，
\eventanchor{TK-DENS-0239-01}{幼帝即位的双辅政格局}由此成形。次年曹爽的亲近者
\eventanchor{TK-DENS-0240-01}{进入尚书、禁军和中书周边}，司马懿转居太傅；
到正始七年，前者继续占据要职，后者以病老姿态减少公开政务，即
\eventanchor{TK-DENS-0246-01}{两套近畿网络暂时分开运作}。帝纪和通鉴可追踪
任官次序，无法还原每次朝议的完整立场；可见的是，调度诏令和宫卫的入口已不再
只是一项辅政名义。

\subsection{从辅政到政权入口}

\eventanchor{TK-0249-GAOPING-TOMBS}{高平陵政变}显示，控制洛阳的宫门和军队可以
把“辅政”转为实际的裁决权。此后王凌、毌丘俭与文钦、诸葛诞在淮南的反抗，
不是同一种动机的重复，却都说明地方军镇会对中枢控制权作出判断。寿春最终在
\eventanchor{TK-0258-ZHUGE-DAN-DEFEAT}{二五八年陷落}，使司马昭清除了最危险的
军事挑战，却也让曹魏皇帝的名义更加空薄。

政变之后，曹爽、何晏等被处死，
\eventanchor{TK-DENS-0249-01}{清洗使司马懿重新控制洛阳朝政}。刑罚的名单与
后来的政权转移相连，却不能把纪传所载的罪名直接等同于全部政治关系的原貌。

废曹芳与
\eventanchor{TK-0260-CAO-MAO-KILLED}{曹髦被杀}是同一条长线的两个刻度：皇帝仍
拥有礼仪位置，却难以独自调动保卫这一位置的力量。司马氏随后受晋公、晋王之
封，终于在\eventanchor{TK-0265-WEI-ABDICATION}{魏晋禅代}中取得名义。这不是
“魏突然消失”，而是多年中军、辅政、封国和近畿控制累积后的结果。

高平陵之后的清除使这一入口越来越窄：王凌案牵出另立宗室的可能，司马懿令
\eventanchor{TK-DENS-0251-01}{楚王曹彪自尽}；正元元年，李丰、夏侯玄等以谋反罪
被杀，\eventanchor{TK-DENS-0254-01}{案情的细节主要来自胜者控制的叙事}，不宜
替它下无保留的定性。甘露元年司马昭继任大将军，
\eventanchor{TK-DENS-0256-01}{整顿中军和淮南防务}，把此前的近畿优势接到边防。
曹髦在南阙附近遇杀的基本顺序可核，随后的
\eventanchor{TK-DENS-0260-01}{近侍出宫行动}及其动机却不能由纪传写成完整的宫廷
现场。

司马昭在次年\eventanchor{TK-DENS-0261-01}{立曹奂而仍以魏廷发令}，并未立即受禅；
咸熙元年再\eventanchor{TK-DENS-0264-01}{进封晋王}，才使封国与礼仪框架更为完备。
这两个刻度恰好留下名义仍在、军政已经转手的间隙。它们既不能证明曹魏在二六四年
已结束，也不能把二六五年的禅代解释为单日完成的夺取。

\subsection{制度得失}

洛阳中枢解决了北方大规模调度的协调问题，也把安全高度集中在宫廷、禁军和
少数辅政者之间。它能迅速组织辽东远征或淮南镇压，却使继承危机直接成为国家
危机。魏的制度得失因而不是“中央集权”这一抽象词，而是近距离的军政资源在
幼主和辅政之间如何被反复争夺。

\begin{sourcenote}[读魏书的尺度]
《三国志》以魏纪年组织全书，最适于追踪朝廷年序，却会影响蜀、吴事件在叙事
中的位置。对高平陵、淮南和禅代，应同时参看帝纪、列传与《资治通鉴》；它们
可互证顺序，不能把参与者的动机化为无争议的事实。
\end{sourcenote}
